% Task 3: the text of the task. Tasks.tex puts all tasks into one PDF; the code shown here is in Task3/.
\settitle{Task 3}{Statistics in a 5G base station}

\begin{story}
A 5G base station receives millions of packets per second from the phones around it, and eight threads
process them. The station keeps three statistics: how many packets have come so far, how many bytes they
had together (for monitoring and billing), and the length of the longest packet so far (to size the buffers).
The code is correct, but the profiler shows that the threads spend most of their time waiting for one mutex.
Make it faster.
\end{story}

Note: this is a very simplified model, made for teaching.

\section{The situation}

\begin{center}
\begin{tikzpicture}[
    pk/.style={draw=nLine, fill=nFill, minimum width=0.36cm, minimum height=0.22cm, inner sep=0pt,
      rounded corners=1pt},
    lane/.style={nLight, thick},
    expl/.style={anchor=north west, text width=7.2cm, font=\small, text=nInk, align=left,
      execute at begin node={\hyphenpenalty=10000\relax}},
    stage/.style={font=\small\bfseries, text=nInk, anchor=west}]
  \def\cx{3.15}
  % ======================= 1. loose =======================
  \foreach \i in {0,...,7} {
    \ifnum\i=3
      \node[actor=sIndigo, minimum width=0.6cm, inner sep=2pt, font=\scriptsize\bfseries] at (\i*0.9,0) {T\i};
    \else
      \node[suspended=sIndigo, minimum width=0.6cm, inner sep=2pt, font=\scriptsize\bfseries] at (\i*0.9,0) {T\i};
    \fi
    \draw[lane] (\i*0.9,-0.3) -- (\i*0.9,-3.0) -- (\cx+\i*0.7-3.5*0.7,-4.2);
  }
  \foreach \i/\y in {0/-0.8, 0/-2.3, 1/-1.4, 1/-2.8, 2/-1.0, 2/-2.0, 3/-0.7, 3/-2.5, 4/-1.7,
                     5/-1.1, 5/-2.4, 6/-0.6, 6/-1.9, 7/-1.3, 7/-2.7}
    \node[pk] at (\i*0.9,\y) {};
  % ======================= 2. the jam =======================
  \fill[AccentOrange!12, rounded corners=4pt] (0.45,-4.15) -- (5.85,-4.15) -- (4.15,-5.75) -- (2.15,-5.75) -- cycle;
  \foreach \r/\n in {0/12, 1/10, 2/8, 3/6, 4/4} {
    \pgfmathtruncatemacro{\last}{\n-1}
    \foreach \j in {0,...,\last}
      \node[pk] at (\cx-\n*0.2+0.2+\j*0.4, -4.4-\r*0.29) {};
  }
  \node[font=\scriptsize\bfseries, text=AccentOrange, anchor=west] at (4.6,-5.45) {the jam};
  % ======================= 3. one at a time =======================
  \node[hw, minimum width=6.0cm, minimum height=2.75cm, anchor=north] (mx) at (\cx,-6.05) {};
  \node[font=\small, text=nInk, anchor=north west] at ($(mx.north west)+(0.1,-0.08)$) {\faLock\ \textbf{mutex m}};
  \draw[-{Stealth}, thick, nInk] (\cx,-5.75) -- (\cx,-6.0);
  \node[pk, minimum width=0.75cm, minimum height=0.32cm, font=\scriptsize\ttfamily, text=nInk, fill=white]
    at (\cx+1.85,-6.3) {512};
  \node[font=\scriptsize, text=nInk, anchor=east] at (\cx+1.4,-6.3) {T3 with};
  \node[font=\small, text=nInk, anchor=north] at (\cx,-6.65) {%
    \begin{tabular}{@{}l@{\quad}r@{\,}c@{\,}r@{}}
      & \scriptsize before & & \scriptsize after \\
      \texttt{packet\_count} & \texttt{1\,041} & $\to$ & \texttt{1\,042} \\
      \texttt{total\_bytes} & \texttt{793\,410} & $\to$ & \texttt{793\,922} \\
      \texttt{longest\_packet} & \texttt{1\,463} & $\to$ & \texttt{1\,463}
    \end{tabular}};
  % ======================= 4. loose again =======================
  \draw[lane] (\cx,-8.8) -- (\cx,-9.1);
  \foreach \i in {0,...,7}
    \draw[lane, -{Stealth}] (\cx,-9.1) -- (\i*0.9,-9.9) -- (\i*0.9,-11.6);
  \foreach \i/\y in {0/-10.3, 1/-11.0, 2/-10.5, 3/-10.1, 3/-11.2, 4/-10.7, 5/-10.2, 6/-10.9, 7/-10.4}
    \node[pk] at (\i*0.9,\y) {};
  % ======================= the explanation on the right =======================
  \begin{scope}[xshift=7.6cm]
    \node[stage] at (0,-0.5) {\badge{1}\ Loose};
    \node[expl] at (0,-0.85) {Eight threads, each with its own packets. Every thread takes a packet and calls
      \texttt{on\_packet(len)}. Plenty of room: the threads work at the same time.};
    \node[stage] at (0,-4.2) {\badge{2}\ The jam};
    \node[expl] at (0,-4.55) {All of them want the same mutex. Only one can go in; the other seven wait in
      front of the door. More threads only make the jam longer.};
    \node[stage] at (0,-6.4) {\badge{3}\ One at a time};
    \node[expl] at (0,-6.75) {Inside the mutex, T3 updates the statistics with its packet of 512 bytes,
      and gets the number \texttt{1\,041} back.};
    \node[stage] at (0,-10.0) {\badge{4}\ Loose again};
    \node[expl] at (0,-10.35) {After the mutex every thread goes on alone with its packet: no more waiting.};
  \end{scope}
  \node[font=\small\bfseries, text=AccentOrange] at (7.9,-12.2)
    {8 threads, but at the mutex only 1 of them works at a time};
\end{tikzpicture}
\end{center}

\codesection

The template is in \ghfolder{Task3}. This is \texttt{stats.cpp}:

\codefile[firstline=1]{stats.cpp}

\runline

The program measures the speed: 8 threads send 1\,000\,000 packets each, and it prints how many million
packets per second \texttt{on\_packet()} handled.

\section{Your tasks}

\begin{questions}
  \item Run it. How many packets per second does the version with the mutex handle?
  \item Replace the mutex with \texttt{std::atomic}. To check your solution, change \texttt{TEST} from 0 to 1 at the
        top of \texttt{stats.cpp}: the program then does not measure, but runs ready tests. Do all tests pass?
        If not, why not, when every variable is atomic?
  \item Set \texttt{TEST} back to 0. How many times faster is it now?
\end{questions}

\begin{note}{The tests (\texttt{TEST 1})}
They are ready in \texttt{packets.h}; you do not write them. 8 threads send packets, and the program checks four
things: the sequence numbers are unique, the byte counter is exact, \texttt{longest\_packet} never goes down, and
\texttt{longest\_packet} is right at the end.
\end{note}

\Needspace{16\baselineskip}
\begin{note}{Atomic variables}
\texttt{std::atomic<int>} (from \texttt{<atomic>}) is an \texttt{int} whose every operation happens in
\textbf{one step}: no other thread can come in the middle of it. In the same way \texttt{std::atomic<uint32\_t>}
is an atomic \texttt{uint32\_t}.
\begin{code}
std::atomic<int> x{5};       // declare: an atomic int that starts at 5

int old = x.fetch_add(3);    // add 3, return the value before: x = 8, old = 5
x++;                         // the same as x.fetch_add(1): x = 9

int now = x.load();          // read: now = 9
if (x > 7) { /* ... */ }     // a plain read works too, it is the same as x.load() > 7

x = 10;                      // write: x = 10
\end{code}
\end{note}
